\chapter*{General Conclusion}
\addcontentsline{toc}{chapter}{General Conclusion}
\markboth{\textbf{General Conclusion}}{}

This report has presented the design and development of a multi-vendor
e-commerce marketplace, carried out as a final-year engineering project
within the host company. The work set out to address concrete gaps in the
local e-commerce landscape --- heavyweight seller onboarding, the absence
of integrated payment, and the lack of any quality moderation --- with a
platform on which independent sellers open a store, publish a catalogue,
and sell to customers under a single, trustworthy roof.

The first two chapters framed the problem and the solution. After
situating the project in its organisational and competitive context, we
analysed the requirements, identified the actors, and specified the
functional and non-functional needs, which we modelled with UML use-case
and class diagrams. The Scrum framework then organised the work into a
product backlog of thirty-six user stories and a plan of two releases of
two sprints each, underpinned by the architecture and technical choices
that frame the realisation.

The last two chapters delivered that plan. \textbf{Release~1} built the
foundations of the marketplace: an identity layer with email
verification, role-based access and database-backed sessions; the
seller-onboarding loop of application and administrative review; and the
core selling capabilities --- store creation, a seven-section WYSIWYG
storefront editor with AI-assisted redesign, product management, and
store-wide sales. \textbf{Release~2} turned the populated catalogue into
a complete product: a multi-store cart and checkout with a
server-authoritative Stripe payment flow confirmed by a
signature-verified webhook; the administration and moderation tooling; an
automated AI verification pipeline with administrative override and a
natural-language database-chat assistant; a premium subscription tier;
and the quality and delivery layers --- a layered suite of 445 automated
tests, containerised services behind a TLS-terminating reverse proxy,
infrastructure described as code with Terraform, and a continuous
deployment pipeline that tests, builds and ships on every push to the
main branch.

Beyond the deliverable itself, the project was an opportunity to practise
the full software-engineering cycle end to end --- from requirements
analysis and agile planning to test-driven implementation, AI
integration and modern DevOps --- and to make pragmatic engineering
decisions under real constraints, such as deriving payment amounts on the
server, screening user content automatically, and keeping the whole
environment reproducible.

Several directions remain open for future work. The features deliberately
left out of scope --- a native mobile application, multi-currency
settlement, advanced fraud detection, an on-premise language model, and a
dedicated full-text search engine --- are natural next steps. The AI
verification and redesign workflows could be extended with richer models
and human-in-the-loop feedback, the analytics could grow into seller- and
marketplace-level dashboards, and the single-droplet deployment could
evolve towards a horizontally scaled, observable production environment.
Together these would carry the platform from a complete, validated
prototype towards a production marketplace at scale.
